% Part: proof-theory
% Chapter: proof-search
% Section: introduction

\documentclass[../../../include/open-logic-section]{subfiles}

\begin{document}

\olfileid{pt}{ps}{int}

\olsection{ভূমিকা}

\Log{G3c} বা \Log{N1i}-র মতো
!!{proof} পদ্ধতির স্বতঃসিদ্ধ
ও অনুমানবিধি নির্ধারণ করে,
সিকোয়েন্ট বা !!{formula}s-এর
কোন বৃক্ষগুলি সঠিক !!{proof}s।
সিকোয়েন্ট বা !!{formula}s-এর
এমন একটি বৃক্ষ দেওয়া থাকলে
(প্রতিটি ধাপে প্রয়োগ করা বিধি,
অনুমানের চিহ্ন ও সেগুলি
অবমুক্তকারী বিধির চিহ্নের
মতো বাড়তি তথ্যসহও হতে পারে)
স্বতঃসিদ্ধ ও বিধির সংজ্ঞা
দিয়ে বৃক্ষটি সঠিক !!{proof}
কি না যাচাই করা যায়।
কিন্তু নির্দিষ্ট সিকোয়েন্টের,
অথবা নির্দিষ্ট অনুমানগুলি থেকে
নির্দিষ্ট !!{formula}-র
!!a{proof} \emph{খুঁজে পাওয়া}
আলাদা এবং সাধারণত আরও
কঠিন সমস্যা। সেটিই
\emph{প্রমাণ-অন্বেষণ}-এর সমস্যা।

প্রমাণ-অন্বেষণের একটি
খুব সরল পদ্ধতি আছে।
!!{formula}s-কে সসীম
বর্ণমালার চিহ্নের অনুক্রম
ভাবতে পারি। সিকোয়েন্ট
এবং সিকোয়েন্ট বা
!!{formula}s-এর বৃক্ষকেও
!!{formula}s-এর অনুক্রম ভাবতে পারি
(শাখা বোঝাতে বিশেষ বন্ধনীসহ
হতে পারে)। স্বতঃসিদ্ধ
ও বিধি দিয়ে কোনো
চিহ্ন-অনুক্রম সঠিক
!!{proof} বোঝায় কি না
যান্ত্রিকভাবে যাচাই করা যায়।
তাই \emph{সম্ভাব্য সব}
সঠিক !!{proof}s-ও
যান্ত্রিকভাবে তৈরি করা যায়:
!!{proof}s বোঝানোর
বর্ণমালায় সব চিহ্ন-অনুক্রম
তৈরি করে প্রতিটি প্রার্থী
সঠিক !!{proof} কি না
যাচাই করতে হয়।
এই সরল প্রমাণ-অন্বেষণ
পদ্ধতিতে সব সঠিক
!!{proof}s তৈরি করতে
থাকি, যতক্ষণ না দেওয়া
সিকোয়েন্টের একটি
!!{proof} পাই (অথবা
দেওয়া অনুমানগুলির মধ্যে
অবমুক্ত না হওয়া অনুমান
রেখে দেওয়া !!{formula}-র
প্রমাণ পাই)।

আরও ভালো পদ্ধতি হলো হাতে
!!a{proof} খোঁজার স্বাভাবিক
রীতিটি নেওয়া: শেষ-সিকোয়েন্ট
থেকে শুরু করে !!a{formula}
বাছি এবং একটি অনুমানবিধি
``উল্টো দিকে'' প্রয়োগ
করে এক বা একাধিক পূর্বধারণা
তৈরি করি। সেগুলির
!!{proof}s থাকলে শেষ
বিধিপ্রয়োগের সঙ্গে মিলিয়ে
দেওয়া সিকোয়েন্টের
!!a{proof} পাওয়া যেত।
স্বাভাবিক নিষ্পাদন পদ্ধতিতে
!!a{proof} খুঁজতেও
একই ধরনের, তবে আরও জটিল,
রীতি ব্যবহার করা যায়।
% BN-SRC-891: translation-local natural-deduction terminology corrected.

কিন্তু এই পদ্ধতি নিশ্চিতভাবে
ফলদায়ক নয়: ভুল বিধি
বা ভুল ক্রমে !!{formula}s
বাছলে আর এগোনোর পথ
নাও থাকতে পারে। পদ্ধতিটি
পুরো নির্দিষ্টও নয়:
প্রতিটি ধাপে বাছার জন্য
একাধিক !!{formula} থাকতে পারে,
বা উল্টো দিকে প্রয়োগের
জন্য একাধিক বিধি থাকতে পারে।
\emph{প্রমাণ-অন্বেষণ অ্যালগরিদম}
হলো সম্পূর্ণ নির্দিষ্ট এমন
পদ্ধতি, যা !!a{proof}
থাকলে তা খুঁজে পাবে
(অর্থাৎ নিশ্চিতভাবে ফলদায়ক)।

কিছু !!{proof} পদ্ধতির
প্রমাণ-অন্বেষণ অ্যালগরিদম
অন্যগুলির চেয়ে সরল হয়।
সিকোয়েন্ট পদ্ধতিতে
!!a{formula}-র !!{main operator}
এবং সূত্রটি বাঁ না ডান দিকে
আছে, তা উল্টো দিকে কোন
বিধি প্রয়োগ করতে হবে
নির্ধারণ করে। যেমন,
$\Gamma \Sequent \Delta, !A \land !B$-র
!!a{proof} খুঁজতে চাই,
যেখানে $!A \land !B$
ডান দিকে আছে। তখন
\RightR{\land} উল্টো দিকে
প্রয়োগ করে সংশ্লিষ্ট দুই
পূর্বধারণা
$\Gamma \Sequent \Delta, !A$
ও $\Gamma \Sequent \Delta, !B$
তৈরি করা উচিত।
কিন্তু \Log{G1c} পদ্ধতিতে
সংকোচন বিধি থাকায়
বিকল্প বেড়ে যায়।
উল্টো দিকে
\RightR{\Contraction}
প্রয়োগ করেও
$\Gamma \Sequent \Delta, !A \land
!B, !A \land !B$
পূর্বধারণা তৈরি করতে পারি।
তারপর আবার
$\Gamma \Sequent
\Delta, !A \land !B, !A \land !B, !A \land !B$,
এভাবে চলতে পারে।
\Log{G2c}-তে সমস্যা আরও বেশি।
\RightR{\land} উল্টো দিকে
প্রয়োগের জন্যও
$\Gamma_1$, $\Gamma_2$
অনুমান করতে হয়,
যাতে $\Gamma
= \Gamma_1 \cup \Gamma_2$;
একইভাবে $\Delta_1$, $\Delta_2$
অনুমান করতে হয়, যাতে
$\Delta
= \Delta_1 \cup \Delta_2$।
তারপর
$\Gamma_1 \Sequent \Delta_1, !A$
ও $\Gamma_2 \Sequent \Delta_2, !B$
পূর্বধারণা নিয়ে চেষ্টা করি।
ভুল অনুমানে পরে আর
এগোনোর পথ না থাকলে
শুরুতে ফিরে অন্য
$\Gamma_1$, $\Gamma_2$,
$\Delta_1$, $\Delta_2$
বাছতে হয়। আবার
!!{formula} যদি
$!A \lor !B$ হয়,
\RightR{\lor}-এর দুই
সম্ভাবনার একটি বেছে
$\Gamma \Sequent \Delta,
!A$ অথবা
$\Gamma \Sequent \Delta, !B$
পূর্বধারণা তৈরি করতে হয়।
এখানেও ভুল পছন্দে
পিছিয়ে ফিরে আসতে হয়।

\Log{G3c}-তে এই সমস্যাগুলি
নেই। তাতে গঠনগত বিধি
নেই; !!{main operator}
এবং !!{formula} কোন
দিকে আছে, তা দিয়েই
অনুমানবিধি বাছা যায়।
তাই \Log{G1c} বা
\Log{G2c}-র তুলনায়
\Log{G3c} প্রমাণ-অন্বেষণের
পক্ষে বেশি উপযোগী।
সফল অন্বেষণে \Log{G3c}-তে
!!a{proof} মেলে; সেই
!!{proof} পরে \Log{G1c}
বা \Log{G2c}-তে
(এমনকি \Log{N1c}-তেও)
অনুবাদ করা যায়।
সুতরাং \Log{G3c}-র জন্য
প্রমাণ-অন্বেষণ অ্যালগরিদম
এবং তার !!{proof}s-কে
অন্য পদ্ধতির !!{proof}s-এ
অনুবাদের অ্যালগরিদম থাকলে
সেই পদ্ধতিগুলির জন্যও
প্রমাণ-অন্বেষণের সমস্যা
সমাধান হয়।

প্রমাণ-অন্বেষণ অ্যালগরিদম
নিশ্চিতভাবে ফলদায়ক কি না
কীভাবে জানব? দেখাতে হবে,
সেটি !!a{proof}
খুঁজে না পেলে কোনো
!!{proof}-ই থাকতে পারে না।
খারাপ অ্যালগরিদম বাছলে
প্রমাণ থাকা সত্ত্বেও
সেটি খুঁজে না পাওয়ার
সম্ভাবনা থাকে।
আগের সরল পদ্ধতিতে
এই সমস্যা নেই:
সেটি সম্ভাব্য সব
!!{proof}s খোঁজে।
কিন্তু প্রমাণ-অন্বেষণ
অ্যালগরিদমের উদ্দেশ্য
দক্ষভাবে যত কম সম্ভব
প্রচেষ্টা খরচ করা।

নির্ভরযোগ্যতা দেখানোর
মূল কৌশল হলো,
অ্যালগরিদম ব্যর্থ হলে
ব্যর্থতার ধরন থেকেই
সিকোয়েন্টটির
!!a{proof} অসম্ভব
দেখানো। এ জন্য
সিকোয়েন্টটি অবৈধ
দেখাই। ধ্রুপদি প্রথম-ক্রমের
যুক্তিতে
$\Gamma \Sequent \Delta$
বৈধ মানে প্রত্যেক
!!{structure}-তে
$\Gamma$-র অন্তত একটি
!!{formula} মিথ্যা,
অথবা~$\Delta$-র অন্তত
একটি !!{formula} সত্য।
\Log{G3c} \emph{যথার্থ}:
এতে শুধু বৈধ
সিকোয়েন্টই !!{prove}s হয়।
!!{proof}s-এর উপর
আরোহ করে তা প্রতিষ্ঠা করা যায়:
স্বতঃসিদ্ধগুলি স্পষ্টতই বৈধ;
বিধির পূর্বধারণাগুলি
বৈধ হলে সিদ্ধান্তও বৈধ।
প্রমাণ-অন্বেষণ অ্যালগরিদমের
নির্ভরযোগ্যতা দেখাতে
প্রমাণ করি, যদি
$\Gamma \Sequent \Delta$-র
!!a{proof} খোঁজা
ব্যর্থ হয়, তাহলে
এমন !!a{structure}
গড়া যায় যাতে
$\Gamma$-র সব
!!{formula}s সত্য
আর~$\Delta$-র সব
!!{formula}s মিথ্যা।
এতে \Log{G3c}-র
\emph{সম্পূর্ণতা}ও মেলে:
$\Gamma \Sequent \Delta$
বৈধ হলে তার
\Log{G3c}-!!{proof} আছে।
প্রতিটি ব্যর্থ অন্বেষণ
$\Gamma \Sequent \Delta$-কে
অবৈধ দেখায়; তাই
বৈধ সিকোয়েন্টের জন্য
প্রমাণ-অন্বেষণ অবশ্যই
সফল হবে।

\end{document}
